\section*{Chapter \ref{ch:b1:weathering-clays} --- Weathering and Clays}

\begin{solution}{exo:b1:weathering-clays:1}
\ce{K+}: $1/1.38 = 0.72$; \ce{Fe^2+}: $2/0.78 = 2.56$; both below 3, soluble
cations. \ce{Fe^3+}: $3/0.645 = 4.65$; \ce{Ti^4+}: $4/0.605 = 6.61$; both
between 3 and 10, hydrolysates. In a waterlogged soil, without oxygen, iron is
released as \ce{Fe^2+}, which stays in solution and is carried away by the
water; in a well-drained soil, oxygen turns it into \ce{Fe^3+}, which
precipitates at once as oxides and hydroxides and stays.
\end{solution}

\begin{solution}{exo:b1:weathering-clays:2}
A swells with glycol and collapses on heating: a \omterm{def:b1:weathering-clays:clay-mineral}{smectite}. B is at
\SI{7}{\angstrom} and is destroyed at \SI{550}{\celsius}: \omterm{def:b1:weathering-clays:clay-mineral}{kaolinite}. C, at
\SI{10}{\angstrom} whatever the treatment: \omterm{def:b1:weathering-clays:clay-mineral}{illite}. D, at \SI{14}{\angstrom}
without swelling or collapse: \omterm{def:b1:weathering-clays:clay-mineral}{chlorite}.
\end{solution}

\begin{solution}{exo:b1:weathering-clays:3}
Two orthoclase bring the 2 aluminium of one \omterm{def:b1:weathering-clays:clay-mineral}{kaolinite}; $6 - 2 = 4$ silicon
leave as \ce{H4SiO4}; 2 \ce{K+} leave, balanced by 2 \ce{H+}; oxygen,
$16 + w = 9 + 16$, gives $w = 9$:
\[
  \ce{2 KAlSi3O8 + 2 H+ + 9 H2O -> Al2Si2O5(OH)4 + 2 K+ + 4 H4SiO4}.
\]
Hydrogen checks: $2 + 18 = 20 = 4 + 16$. All the potassium leaves, and $4/6 =
2/3$ of the silicon.
\end{solution}

\begin{solution}{exo:b1:weathering-clays:4}
(a) The mass does not change: $10 \times 999.87/916.72 = \qty{10.91}{cm^3}$.
(b) Between about \qty{-4}{\celsius} and \qty{-15}{\celsius}. Just below zero
the ice cannot build up enough pressure to wedge the crack; in a harder frost
the water no longer moves through the pores toward the ice, which stops
growing. (c) By
\cref{prop:b1:weathering-clays:surface} with $n = \qty{1}{m}/\qty{1}{cm} = 100$,
the area is multiplied by 100: from \qty{6}{m^2} to \qty{600}{m^2}.
\end{solution}

\begin{solution}{exo:b1:weathering-clays:5}
$S/m = 6/(\rho d)$: for \qty{1}{cm}, $6/(2600\times0.01) =
\qty{0.231}{m^2.kg^{-1}}$; for \qty{1}{mm}, \qty{2.31}{m^2.kg^{-1}}; for
\qty{1}{\micro m}, \qty{2308}{m^2.kg^{-1}}, or \qty{2.3}{m^2.g^{-1}}.
Chemical reactions take place at the surface of the grains, so their rate per
kilogram of rock grows with the specific surface, ten thousand times from the
centimetre to the micrometre. \omterm{def:b1:weathering-clays:weathering}{Physical weathering} cuts the rock into ever
smaller pieces and opens cracks through which water enters: it multiplies the
surface on which \omterm{def:b1:weathering-clays:weathering}{chemical weathering} can act.
\end{solution}

\begin{solution}{exo:b1:weathering-clays:6}
A tonne of albite is $\num{e6}/262.22 = \qty{3813.6}{mol}$. By
\cref{prop:b1:weathering-clays:feldspars}, two moles give one of \omterm{def:b1:weathering-clays:clay-mineral}{kaolinite}:
$\qty{1906.8}{mol} \times \qty{258.16}{g.mol^{-1}} = \qty{492.3}{kg}$. Two of
the three silicon atoms leave: $2\times3813.6 = \qty{7627.2}{mol}$ of
\ce{SiO2}, \qty{458.2}{kg}. All the sodium leaves: $3813.6\times22.99 =
\qty{87.7}{kg}$. Less than half the mass of the \omterm{def:b1:rock-forming-minerals:groups}{feldspar} remains as clay.
\end{solution}

\begin{solution}{exo:b1:weathering-clays:7}
By \cref{prop:b1:weathering-clays:immobile}, a kilogram of granite leaves
$0.30/0.60 = \qty{0.5}{kg}$ of saprolite. Fraction lost: sodium $1 -
(0.35/3.5)(0.5) = 0.95$; silicon $1 - (60/72)(0.5) = 0.583$; aluminium $1 -
(27/14)(0.5) = 0.036$. Aluminium is practically kept, the sodium has almost
entirely gone, and nearly \qty{60}{\percent} of the silicon too, although much of what
remains is held in quartz: the \omterm{def:b1:rock-forming-minerals:groups}{feldspars} are being turned into \omterm{def:b1:weathering-clays:clay-mineral}{kaolinite},
\omterm{def:b1:weathering-clays:hydrolysis-degree}{monosiallitisation}.
\end{solution}

\begin{solution}{exo:b1:weathering-clays:8}
Pyrophyllite: $x = y = 0$, charge 0. Muscovite layer: $x = 1$, charge $-1$.
Beidellite layer: $x = 0.5$, charge $-0.5$. Montmorillonite layer: $y = 0.33$,
charge $-0.33$. The two \omterm{def:b1:weathering-clays:clay-mineral}{smectite} layers, beidellite and montmorillonite, can
swell: their small charge is balanced by few interlayer cations, which hold the
layers loosely and draw water in as they hydrate. Pyrophyllite has no
interlayer cations to hydrate, and the charge $-1$ of muscovite is balanced by
potassium ions that lock the layers.
\end{solution}

\begin{solution}{exo:b1:weathering-clays:9}
(a) Mg: 2 and 2; Si: 1 and 1; C: 4 and 4; O: $4 + 8 + 4 = 16$ and $12 + 4 =
16$; H: 8 and $4 + 4 = 8$; charge: 0 and $4 - 4 = 0$. (b) The molar mass of
forsterite is $2\times24.305 + 28.085 + 4\times15.999 =
\qty{140.69}{g.mol^{-1}}$, so a tonne is \qty{7108}{mol}. It consumes
$4\times7108 = \qty{28431}{mol}$ of \ce{CO2}, \qty{1251}{kg}; with two
magnesium per formula unit, \cref{prop:b1:weathering-clays:co2} stores
$2\times7108$ moles, \qty{626}{kg}, half the amount consumed. (c) A tonne of
calcite, \qty{9991}{mol}, consumes \qty{9991}{mol} of \ce{CO2} (\qty{440}{kg})
on dissolving, and gives all of it back when the carbonate precipitates again:
nothing is stored.
\end{solution}

\begin{solution}{exo:b1:weathering-clays:10}
(a) $3 < 3/r < 10$ gives $\qty{0.3}{\angstrom} < r < \qty{1.0}{\angstrom}$.
(b) $3/0.27 = 11.1 > 10$: boron travels as an oxyanion, borate. (c) $2/0.83 =
2.41$, a soluble cation; $4/0.53 = 7.55$, a hydrolysate. In water poor in
oxygen manganese is dissolved as \ce{Mn^2+}; where the spring meets the air,
it is oxidised to \ce{Mn^4+}, which precipitates at once as a black oxide on
the rock. (d) Its \omterm{def:b1:weathering-clays:ionic-potential}{ionic potential}, $4/0.40 = 10.0$, lies on the boundary
between hydrolysates and oxyanions: silicon is neither held like aluminium nor
carried off like sulfur, and the fraction leached depends on how much water
passes, the basis of the three degrees of \omterm{def:b1:weathering-clays:hydrolysis}{hydrolysis}.
\end{solution}

\begin{solution}{exo:b1:weathering-clays:11}
(a) With \qty{12.011}{g.mol^{-1}} of carbon, \num{11.7e12} to
\qty{13.6e12}{mol} is 0.141 to 0.163 billion tonnes of carbon a year; with
\qty{44.009}{g.mol^{-1}}, 0.515 to 0.599 billion tonnes of carbon dioxide.
(b) $4.08/13.6 = \qty{30}{\percent}$ to $4.08/11.7 = \qty{35}{\percent}$.
(c) Half the carbon dioxide consumed is stored when the calcium and magnesium
precipitate as carbonate: 0.070 to 0.082 billion tonnes of carbon a year. The
real figure is smaller, because part of the bicarbonate goes with sodium and
potassium, which form no carbonate on the sea floor.
\end{solution}

\begin{solution}{exo:b1:weathering-clays:12}
Si/Al $= (w_{\ce{SiO2}}/60.08)/(2w_{\ce{Al2O3}}/101.96)$. A: $0.7748/0.7748 =
1.00$, the ratio of \omterm{def:b1:weathering-clays:clay-mineral}{kaolinite}, \omterm{def:b1:weathering-clays:hydrolysis-degree}{monosiallitisation}. B: $0.9987/0.4904 = 2.04$, a
2:1 \omterm{def:b1:weathering-clays:clay-mineral}{clay mineral}, \omterm{def:b1:weathering-clays:hydrolysis-degree}{bisiallitisation}. C: $0.0333/1.2358 = 0.027$, nearly no
silicon left, gibbsite, \omterm{def:b1:weathering-clays:hydrolysis-degree}{allitisation}. From the least to the most intense
leaching, that is toward the hottest and wettest climate: B, A, C.
\end{solution}

\begin{solution}{pb:b1:weathering-clays:1}
\textbf{1.} \ce{2 NaAlSi3O8 + 2 H+ + 9 H2O -> Al2Si2O5(OH)4 + 2 Na+ + 4 H4SiO4}.

\textbf{2.} \ce{CaAl2Si2O8 + 2 H+ + H2O -> Al2Si2O5(OH)4 + Ca^2+}. No silicon
leaves: anorthite already has the one silicon per aluminium of \omterm{def:b1:weathering-clays:clay-mineral}{kaolinite}.

\textbf{3.} One gibbsite takes one aluminium, half an anorthite: doubling,
\ce{CaAl2Si2O8 + 2 H+ + 6 H2O -> 2 Al(OH)3 + Ca^2+ + 2 H4SiO4}. Oxygen:
$8 + 6 = 6 + 8$; hydrogen: $2 + 12 = 6 + 8$.

\textbf{4.} \ce{5 NaAlSi3O8 + 4 H+ + 16 H2O -> 2 Na_{0.5}Al2(Si_{3.5}Al_{0.5})O10(OH)2 + 4 Na+ + 8 H4SiO4}; % ce-unbalanced-ok: decimal subscripts the checker cannot read; balance tested in tests/bachelor-1/test_weathering_clays.py
hydrogen, $4 + 32 = 36$ on the left and $2\times2 + 8\times4 = 36$ on the
right, balances.

\textbf{5.} Each \ce{H+} comes with one \ce{HCO3-} from one \ce{CO2} and one
\ce{H2O}:
\ce{2 NaAlSi3O8 + 2 CO2 + 11 H2O -> Al2Si2O5(OH)4 + 2 Na+ + 2 HCO3- + 4 H4SiO4}
and \ce{CaAl2Si2O8 + 2 CO2 + 3 H2O -> Al2Si2O5(OH)4 + Ca^2+ + 2 HCO3-}.

\textbf{6.} Albite to beidellite is \omterm{def:b1:weathering-clays:hydrolysis-degree}{bisiallitisation}; albite and anorthite to
\omterm{def:b1:weathering-clays:clay-mineral}{kaolinite} are \omterm{def:b1:weathering-clays:hydrolysis-degree}{monosiallitisation}; anorthite to gibbsite is \omterm{def:b1:weathering-clays:hydrolysis-degree}{allitisation}.

\textbf{7.} $\num{e6}/262.22 = \qty{3813.6}{mol}$ of albite;
$\num{e6}/278.20 = \qty{3594.5}{mol}$ of anorthite.

\textbf{8.} Albite: $3813.6/2 = \qty{1906.8}{mol}$ of \omterm{def:b1:weathering-clays:clay-mineral}{kaolinite},
\qty{492.3}{kg}. Anorthite: \qty{3594.5}{mol}, \qty{928.0}{kg}.

\textbf{9.} Albite: $2\times3813.6 = \qty{7627.2}{mol}$ of \ce{SiO2},
\qty{458.2}{kg}, and \qty{3813.6}{mol} of sodium, \qty{87.7}{kg}. Anorthite:
no silica, and \qty{3594.5}{mol} of calcium, \qty{144.1}{kg}.

\textbf{10.} Albite: $1/2.62 = \qty{0.382}{m^3}$; its \omterm{def:b1:weathering-clays:clay-mineral}{kaolinite},
$0.4923/2.6 = \qty{0.189}{m^3}$: the solid loses \qty{50.4}{\percent} of its
volume. Anorthite: $1/2.75 = \qty{0.364}{m^3}$; its \omterm{def:b1:weathering-clays:clay-mineral}{kaolinite},
$0.9280/2.6 = \qty{0.357}{m^3}$: a loss of only \qty{1.8}{\percent}.

\textbf{11.} The space left becomes pores, filled with water and air: a
saprolite on an albite-rich granite is light, porous and crumbly, its quartz
grains loose in a little clay. Anorthite turns into nearly its own volume of
\omterm{def:b1:weathering-clays:clay-mineral}{kaolinite}, so the weathered gabbro keeps most of its solid matter.

\textbf{12.} Sodium ($z/r = 0.98$) and calcium (2.0) are soluble cations and
leave; aluminium (5.6) is a hydrolysate and stays entirely in the \omterm{def:b1:weathering-clays:clay-mineral}{kaolinite};
silicon (10.0), on the boundary, leaves in part from albite and not at all
from anorthite, whose silicon \omterm{def:b1:weathering-clays:clay-mineral}{kaolinite} can hold.

\textbf{13.} One \ce{CO2} per albite: \qty{3813.6}{mol}, $3813.6 \times 44.01 =
\qty{167.8}{kg}$.

\textbf{14.} Two per anorthite: \qty{7189.0}{mol}, \qty{316.4}{kg}.

\textbf{15.} $7189.0 \times 61.02 = \qty{438.7}{kg}$ of bicarbonate.

\textbf{16.} Organisms precipitate it as calcium carbonate,
\ce{Ca^2+ + 2 HCO3- -> CaCO3 + CO2 + H2O}: for each calcium, one of the two
\ce{CO2} goes back to the air, \qty{3594.5}{mol} or \qty{158.2}{kg}, and the
other is buried in the carbonate.

\textbf{17.} Sodium forms no carbonate on the sea floor: its bicarbonate
cannot be stored as carbonate rock, and the long-term storage is counted from
calcium and magnesium only.

\textbf{18.} A tonne of calcite consumes \qty{9991}{mol} of \ce{CO2} when it
dissolves and gives it all back when it precipitates: nothing is stored, and
only the weathering of calcium and magnesium silicates draws carbon dioxide
down for good.

\textbf{19.} $\qty{200}{m^3.s^{-1}} \times \qty{0.3}{mol.m^{-3}} =
\qty{60}{mol.s^{-1}}$; over a year of \qty{3.156e7}{s}, \qty{1.893e9}{mol}, or
$\num{1.893e9}/\num{50000} = \qty{37870}{mol}$ per square kilometre.

\textbf{20.} Two bicarbonates per anorthite: \qty{9.467e8}{mol} a year,
$\num{9.467e8}\times\qty{278.20}{g} = \qty{263400}{t}$ a year, or
\qty{5.27}{t} per square kilometre.

\textbf{21.} The fresh rock consumed is $5.27/0.20 = \qty{26.3}{t}$, or
$\num{26340}/2700 = \qty{9.75}{m^3}$, per square kilometre and year: a layer
\qty{9.75}{\micro m} thick a year, \qty{9.8}{m} per million years.

\textbf{22.} One \ce{CO2} stored per calcium: $\num{9.467e8} \times 44.01 =
\qty{4.17e10}{g}$, about \qty{41700}{t} a year.

\textbf{23.} It counts only the matter carried away in solution. The \omterm{def:b1:weathering-clays:clay-mineral}{kaolinite}
and quartz left by weathering are also removed by erosion, as mud and sand in
the river, and they are not counted; nor is the weathering of the other
\omterm{def:b1:rock-forming-minerals:mineral}{minerals} of the rock. And if part of the bicarbonate comes from sodium or
potassium \omterm{def:b1:rock-forming-minerals:groups}{feldspars}, which give one bicarbonate per formula unit instead of
two, more rock had to weather to supply it.

\textbf{24.} By \cref{prop:b1:weathering-clays:co2}, one mole of \ce{CO2} is
stored per mole of anorthite: $3594.5 \times \qty{44.01}{g} = \qty{158.19}{kg}$.
The weathering of one tonne of anorthite draws down for good
\textbf{158}~kg of carbon dioxide.
\end{solution}
